\documentclass[11pt]{article}
\usepackage[margin=1in]{geometry}
\usepackage{amsmath,amssymb}
\usepackage{booktabs}
\usepackage{siunitx}
\usepackage{enumitem}
\usepackage{xcolor}
\usepackage[colorlinks=true,linkcolor=blue]{hyperref}

\newcommand{\Ven}{V_{\mathrm{EN}}}
\newcommand{\Vref}{V_{\mathrm{REF}}}
\newcommand{\Vtlv}{V_{\mathrm{TLV}}}
\newcommand{\Vth}{V_{\mathrm{EN,th}}}
\newcommand{\Iref}{I_{\mathrm{ref}}}
\newcommand{\Istby}{I_{\mathrm{STBY}}}
\newcommand{\Ihys}{I_{\mathrm{HYS}}}
\newcommand{\IH}{I_{H}}
\newcommand{\Rone}{R_{1a}}
\newcommand{\Rtwo}{R_{1b}}
\newcommand{\Rt}{R_{2}}

\title{LM5176 \texttt{EN/UVLO} Threshold with TLV431 Latch:\\ Equations, Ranges, and Design Constraints}
\author{}
\date{}

\begin{document}
\maketitle

\section{Circuit topology and conventions}

A single node, \texttt{EN/UVLO} (voltage $\Ven$), is driven from $V_{IN}$ through
$\Rt$. A two--element lower leg returns to ground: $\Rtwo$ connects \texttt{EN/UVLO}
to the TLV431 reference node (voltage $\Vref$), and $\Rone$ connects that reference
node to ground. The TLV431 \texttt{REF} pin sits on the $\Vref$ node and draws a bias
current $\Iref$ \emph{into} the pin.

\begin{center}
\begin{tabular}{ll}
$V_{IN}\;\longrightarrow\;[\,\Rt\,]\;\longrightarrow\;\Ven$ &  (top resistor) \\
$\Ven\;\longrightarrow\;[\,\Rtwo\,]\;\longrightarrow\;\Vref$ & (mid resistor) \\
$\Vref\;\longrightarrow\;[\,\Rone\,]\;\longrightarrow\;\mathrm{GND}$ & (bottom resistor) \\
\end{tabular}
\end{center}

\paragraph{Current--source convention (LM5176).}
The LM5176 \emph{sources} current \emph{out of} the \texttt{EN/UVLO} pin into the node.
In standby it sources $\Istby$. Once the operating (UVLO) threshold is crossed it
sources an \emph{additional} $\Ihys$, which lifts $\Ven$ and creates the turn-on/turn-off
hysteresis. We write the total injected current as
\[
I_S \;=\; \Istby + \IH,\qquad
\IH=\begin{cases}0 & \text{standby (rising, before turn-on)}\\[2pt]\Ihys & \text{running (after turn-on)}\end{cases}
\]
\textbf{Sign check:} verify against your datasheet revision that both currents are
\emph{sourced from} the pin (raising $\Ven$). If your part \emph{sinks} the hysteresis
current instead, flip the sign of the $I_S$ term below.

\section{Parameter ranges}

\begin{center}
\begin{tabular}{lccl}
\toprule
Symbol & Min & Max & Notes \\
\midrule
$V_{IN}$            & \SI{4.5}{\volt}  & \SI{22}{\volt}   & USB input \\
$\Vth$ (operating)  & \SI{1.17}{\volt} & \SI{1.29}{\volt} & LM5176 \texttt{EN/UVLO} operating threshold \\
$\Istby$            & \SI{1}{\micro\ampere}    & \SI{3}{\micro\ampere}    & standby source current \\
$\Ihys$             & \SI{2.15}{\micro\ampere} & \SI{4.25}{\micro\ampere} & added after turn-on \\
$\Iref$             & \SI{0}{\micro\ampere}    & \SI{0.38}{\micro\ampere} & TLV431 \texttt{REF} input current \\
$\Vtlv$             & \SI{1.206}{\volt}& \SI{1.266}{\volt}& effective TLV431 reference (see below) \\
\bottomrule
\end{tabular}
\end{center}

\paragraph{Effective TLV431 reference $\Vtlv$.}
Built up from the nominal \SI{1.24}{\volt}:
\[
\Vtlv^{\max}=1.240+\underbrace{0.0062}_{+0.5\%}+\underbrace{0.020}_{\text{temp}}=\SI{1.266}{\volt},
\]
\[
\Vtlv^{\min}=1.240-\underbrace{0.0062}_{-0.5\%}-\underbrace{0.020}_{\text{temp}}-\underbrace{0.008}_{V_{KA}\text{ effect}}=\SI{1.206}{\volt}.
\]
The $V_{KA}$ term uses $\approx\SI{-1.5}{\milli\volt\per\volt}$ over $V_{KA}$ up to
$\sim\!\SI{6}{\volt}$ (only shifts the reference \emph{down}, so it appears in the min only).

\section{Node equations (KCL)}

\paragraph{\texttt{EN/UVLO} node.}
\begin{equation}
\frac{V_{IN}-\Ven}{\Rt}+I_S \;=\; \frac{\Ven-\Vref}{\Rtwo}
\label{eq:kcl_en}
\end{equation}

\paragraph{\texttt{REF} node.}
\begin{equation}
\frac{\Ven-\Vref}{\Rtwo} \;=\; \frac{\Vref}{\Rone}+\Iref
\label{eq:kcl_ref}
\end{equation}

\section{Closed-form solutions}

Solving \eqref{eq:kcl_en}--\eqref{eq:kcl_ref}:

\begin{equation}
\boxed{\;\Ven \;=\; \frac{V_{IN}(\Rone+\Rtwo)+I_S\,\Rt(\Rone+\Rtwo)-\Iref\,\Rone\,\Rt}{\Rone+\Rtwo+\Rt}\;}
\label{eq:ven}
\end{equation}

with $I_S=\Istby+\IH$. (This matches your stated expression exactly.)

\begin{equation}
\boxed{\;\Vref \;=\; \frac{\Rone\,(\Ven-\Iref\,\Rtwo)}{\Rone+\Rtwo}\;}
\label{eq:vref}
\end{equation}

\paragraph{Correction note.} Equation \eqref{eq:vref} is the form consistent with the
same KCL that produced \eqref{eq:ven}; the bias term is $\Iref\Rtwo$. A previously used
form with $\Iref\Rone$ overstates the correction by $\Iref\,\Rone^{2}/(\Rone+\Rtwo)$
(tens of mV at these values), so \eqref{eq:vref} should be preferred.

\paragraph{TLV431 trip relation (the key design lever).}
The map between $\Ven$ and $\Vref$ depends only on the lower leg, \emph{independent of
$V_{IN}$, $\Rt$, and the EN-pin currents}:
\begin{equation}
\Ven \;=\; \Vref\,\frac{\Rone+\Rtwo}{\Rone}+\Iref\,\Rtwo .
\label{eq:vref_to_ven}
\end{equation}
Setting $\Vref=\Vtlv$ gives the \texttt{EN/UVLO} voltage at the instant the TLV431
crosses its reference:
\begin{equation}
\boxed{\;\Ven\big|_{\mathrm{trip}} \;=\; \Vtlv\,\frac{\Rone+\Rtwo}{\Rone}+\Iref\,\Rtwo\;}
\label{eq:ven_trip}
\end{equation}

\paragraph{Inverting for $V_{IN}$.} For any target $\Ven$ in a given state,
\begin{equation}
V_{IN}\;=\;\frac{\Ven(\Rone+\Rtwo+\Rt)-I_S\,\Rt(\Rone+\Rtwo)+\Iref\,\Rone\,\Rt}{\Rone+\Rtwo}.
\label{eq:vin}
\end{equation}
Useful $V_{IN}$ points:
\begin{itemize}[nosep]
\item \textbf{Turn-on} (rising): $\Ven=\Vth$, $I_S=\Istby$.
\item \textbf{LM5176 natural turn-off} (falling): $\Ven=\Vth$, $I_S=\Istby+\Ihys$.
\item \textbf{TLV431 latch trip} (falling): $\Ven=\Ven|_{\mathrm{trip}}$, $I_S=\Istby+\Ihys$.
\end{itemize}

\section{Latch behavior and operating sequence}

The external latch is powered from $V_{IN}$ and \textbf{defaults to engaged}: it is
held off (\emph{dearmed}) only while $\Vref>\Vtlv$. There is no
``arm-after-first-release'' grace --- whenever $\Vref$ sits below $\Vtlv$ with the
converter running, the latch grabs and pulls \texttt{EN/UVLO} low until a power cycle.
The divider must therefore keep the TLV431 conducting throughout normal operation.

\begin{enumerate}[nosep]
\item $V_{IN}$ rises; $\Ven$ (standby) reaches $\Vth$ at $V_{\mathrm{on}}$, LM5176 turns on.
\item $\IH=\Ihys$ switches in $\Rightarrow$ $\Ven$, hence $\Vref$, jumps up. Because the
design forces $V_{\mathrm{latch}}<V_{\mathrm{on}}$ (constraint~C2), $\Vref$ is already
above $\Vtlv$ at this instant: the latch is dearmed \emph{at} turn-on, leaving no
running-but-not-dearmed window even on an arbitrarily slow rise.
\item $V_{IN}$ falls. While running, $\Vref$ falls; when it drops below $\Vtlv$ at
$V_{\mathrm{latch}}$ the TLV431 stops conducting, the latch engages, \texttt{EN/UVLO} is
pulled low, and it stays latched. This occurs before the LM5176 would reach $\Vth$ on its
own (constraint~C3).
\end{enumerate}

The three running-state events must therefore order, \emph{for every individual unit}, as
\begin{equation}
V_{\mathrm{off}}\;<\;V_{\mathrm{latch}}\;<\;V_{\mathrm{on}},
\label{eq:order}
\end{equation}
i.e.\ the latch trip point lies strictly inside the LM5176's own hysteresis window.

\section{Design constraints (worst-case)}

Each must hold across resistor tolerance ($\pm5\%$) and the parameter extremes noted.

\paragraph{C1 --- Guaranteed turn-on by \SI{4.5}{\volt}.}
Turn-on must be assured before the weakest expected supply (a sagging
\SI{4.5}{\volt} USB rail), so in \emph{standby}
\begin{equation}
\min \;\Ven\big(V_{IN}{=}4.5\;I_S{=}\Istby^{\min},\;\Iref^{\max}\big)\;\ge\;\Vth^{\max}=\SI{1.29}{\volt}.
\end{equation}

\paragraph{C2 --- Latch dearmed below turn-on (relative).}
The latch must be held off at turn-on, which requires $\Vref>\Vtlv$ the instant the
converter starts. Equivalently the latch trip point must sit below this unit's own
turn-on, with margin $\Delta_{\mathrm{hi}}$:
\begin{equation}
\min_{\text{corners}}\big(V_{\mathrm{on}}-V_{\mathrm{latch}}\big)\;\ge\;\Delta_{\mathrm{hi}}
\qquad(\text{e.g.\ }\Delta_{\mathrm{hi}}=\SI{0.1}{\volt}).
\label{eq:c2}
\end{equation}
This is the corrected form of C2. Comparing $V_{\mathrm{latch}}$ to a \emph{fixed}
\SI{4.5}{\volt} would be wrong: on a slow rise the reference that matters is the unit's
own $V_{\mathrm{on}}$, not the spec minimum. If $V_{\mathrm{latch}}>V_{\mathrm{on}}$ there
is a band where the converter runs yet $\Vref<\Vtlv$, and the default-engaged latch grabs.
Like C3, \eqref{eq:c2} is a \emph{per-unit} (correlated) difference --- it must be
evaluated as one subtraction per device, never as the gap between independent population
bands.

\paragraph{C3 --- Latch fires before the LM5176's own UVLO (relative).}
On a falling input the TLV431 must trip before the LM5176 reaches $\Vth$:
\begin{equation}
\min_{\text{corners}}\big(V_{\mathrm{latch}}-V_{\mathrm{off}}\big)\;\ge\;\Delta_{\mathrm{lo}}
\qquad(\text{e.g.\ }\Delta_{\mathrm{lo}}=\SI{0.1}{\volt}).
\label{eq:c3}
\end{equation}

\paragraph{The hysteresis budget links C2 and C3.}
Subtracting the running-state turn-on and dropout expressions gives the LM5176's intrinsic
input hysteresis $V_{\mathrm{on}}-V_{\mathrm{off}}=\Ihys\Rt$, so
\begin{equation}
\boxed{\;\big(V_{\mathrm{on}}-V_{\mathrm{latch}}\big)\;=\;\Ihys\Rt\;-\;\big(V_{\mathrm{latch}}-V_{\mathrm{off}}\big)\;}
\label{eq:budget}
\end{equation}
The two gaps share one budget: $\Delta_{\mathrm{hi}}+\Delta_{\mathrm{lo}}\le\Ihys\Rt$ at the
worst corner. A larger $\Rtwo$ ($\Rt$) widens the window, but the $\Vth$ ($\pm\SI{60}{\milli\volt}$)
and $\Vtlv$ ($\pm\SI{30}{\milli\volt}$) spreads inflate the $V_{\mathrm{latch}}-V_{\mathrm{off}}$
gap by the factor $(\Rone+\Rtwo+\Rt)/(\Rone+\Rtwo)$, so a substantial $\Ihys\Rt$ is needed
just to contain that spread. That forces a large $\Rt$ and hence a low, soft turn-on, with
divider current sinking toward the LM5176 bias currents; the search rejects the degenerate
region (where the bias currents, not $V_{IN}$, set $\Ven$) via a divider/bias current-ratio
guard. The net effect is that \eqref{eq:order} is satisfiable only in a narrow design band
--- a real consequence of the part tolerances, not a modelling artifact.

\paragraph{Worked result.}
A search over E96 values (resistors $\pm5\%$, all parameters at their extremes) gives, as
the highest-turn-on triple still satisfying \eqref{eq:order},
\[
\Rt=\SI{374}{\kilo\ohm},\quad \Rone=\SI{232}{\kilo\ohm},\quad \Rtwo=\SI{27.4}{\kilo\ohm}
\]
(with $\Rone=R_{1a}$, $\Rtwo=R_{1b}$, $\Rt=R_2$): turn-on guaranteed by $\sim\!\SI{3.1}{\volt}$,
$V_{\mathrm{latch}}<V_{\mathrm{on}}$ and $V_{\mathrm{latch}}>V_{\mathrm{off}}$ each by
$\sim\!\SI{0.1}{\volt}$ worst-case. The thin margins and soft turn-on band
(\SIrange{1.6}{3.1}{\volt}) are the cost of \eqref{eq:order} at these tolerances.

\paragraph{Startup caveat.}
Because the latch defaults to engaged, the design must also ensure \texttt{EN/UVLO} can
reach $V_{\mathrm{on}}$ at cold start before the latch can grab (a brief power-up release
or RC delay on the latch network) --- \eqref{eq:order} sizes the steady-state thresholds,
not the power-up race.

\section{Worst-case correlation: why the bands may overlap but units never do}

It is tempting to compute the population span of each $V_{IN}$ event independently and
compare the spans. That is \emph{too pessimistic} for the latch-vs-dropout question.

Both the LM5176 self-dropout ($\Ven=\Vth$) and the TLV431 latch ($\Ven=\Ven|_{\mathrm{trip}}$)
are \textbf{running-state events of the same device at the same instant}. They share the
identical $\Istby$, $\Ihys$, $V_{IN}$, and $\Rt$. Subtracting the two instances of
\eqref{eq:vin} with a common $I_S$ and $\Iref$,
\begin{equation}
\boxed{\;V_{\mathrm{latch}}-V_{\mathrm{off}}
=\big(\Ven|_{\mathrm{trip}}-\Vth\big)\,\frac{\Rone+\Rtwo+\Rt}{\Rone+\Rtwo}\;}
\label{eq:gap}
\end{equation}
and the bracket
\[
\Ven|_{\mathrm{trip}}-\Vth=\Vtlv\,\frac{\Rone+\Rtwo}{\Rone}+\Iref\Rtwo-\Vth
\]
contains \emph{no current term at all}: $\Istby$, $\Ihys$, $V_{IN}$, $\Rt$ have all
cancelled. So the per-unit gap depends only on $\Vtlv$, the $\Rone/\Rtwo$ ratio, $\Iref$,
and $\Vth$. Driving $\Ihys$ to its maximum to lift $V_{\mathrm{off}}$ while
simultaneously driving it to its minimum to drop $V_{\mathrm{latch}}$ describes no
physical unit; it double-counts one parameter.

\paragraph{Consequence.} The independent worst-case bands of $V_{\mathrm{off}}$ and
$V_{\mathrm{latch}}$ can overlap on paper, yet $V_{\mathrm{latch}}-V_{\mathrm{off}}>0$
holds for every individual device as long as $\Ven|_{\mathrm{trip}}>\Vth$ --- this is
constraint~C3, \eqref{eq:c3}. \textbf{The same caution applies to C2.} The corrected C2,
\eqref{eq:c2}, is also a difference of two events on the same device
($V_{\mathrm{on}}-V_{\mathrm{latch}}$): the standby turn-on and the running latch trip,
which share $\Iref$ and the resistors. It too must be differenced per-unit, not read off
the gap between the $V_{\mathrm{on}}$ and $V_{\mathrm{latch}}$ population bands (those bands
overlap heavily while every device still keeps $V_{\mathrm{latch}}<V_{\mathrm{on}}$). The
rule of thumb: compare to a \emph{fixed external} reference $\Rightarrow$ absolute
worst-case is fine; compare two \emph{co-varying internal} quantities $\Rightarrow$ the
difference must be taken per-unit. The earlier ``C2 vs.\ a fixed \SI{4.5}{\volt}'' form
violated this and is superseded. The accompanying script reports both the honest
population bands (for visualisation) and the guaranteed per-unit gaps
\texttt{gap\_lo}~$=\min(V_{\mathrm{latch}}-V_{\mathrm{off}})$ and
\texttt{gap\_hi}~$=\min(V_{\mathrm{on}}-V_{\mathrm{latch}})$, both of which must stay
positive with margin.

\section{Symbol summary}
\begin{center}
\begin{tabular}{ll}
\toprule
$\Rt$ & top resistor, $V_{IN}\to$\texttt{EN/UVLO} \\
$\Rtwo$ & mid resistor, \texttt{EN/UVLO}$\to\Vref$ \\
$\Rone$ & bottom resistor, $\Vref\to$ GND \\
$\Ven$ & \texttt{EN/UVLO} pin voltage \\
$\Vref$ & TLV431 reference-node voltage \\
$\Ven|_{\mathrm{trip}}$ & $\Ven$ when $\Vref=\Vtlv$, Eq.~\eqref{eq:ven_trip} \\
\bottomrule
\end{tabular}
\end{center}

\end{document}
